\section{Introduction}

Graph classification is a fundamental task in graph representation learning, with applications in molecular property prediction, drug discovery, social network analysis, and computer vision. Unlike node-level prediction, graph classification requires the model to summarize an entire graph into a single discriminative representation. This setting amplifies several well-known difficulties of graph neural networks (GNNs): graphs vary in size and topology, local neighborhoods can be noisy, long-range dependencies are hard to preserve, and deep message passing often suffers from over-smoothing and unstable optimization.

Modern GNNs address these challenges through a diverse set of message-passing operators. Graph Convolutional Networks (GCN) \cite{kipf2016semi} emphasize smooth neighborhood aggregation, Graph Attention Networks (GAT) \cite{velivckovic2017graph} learn adaptive neighbor weights, and graph Transformers \cite{dwivedi2020generalization} broaden the receptive field through attention. Despite their differences, these operators are usually studied as alternative backbones inside a largely sequential architecture. As depth increases, repeated aggregation can still blur node representations and weaken the useful diversity of intermediate features \cite{li2018deeper, chen2023revisiting, xu2018powerful}.

Several research directions have tried to alleviate this problem. Residual connections \cite{he2016deep, bresson2017residual}, normalization \cite{zhao2020pair, cai2021graphnorm}, and edge dropout \cite{rong2019dropedge} stabilize deep training. Multi-scale aggregation methods such as Jumping Knowledge \cite{xu2018representation} and dense connectivity \cite{du2024densegnn} preserve information from multiple depths. Pooling methods such as DiffPool \cite{ying2018hierarchical} and SAGPool \cite{lee2019self} add hierarchical structure. However, most of these designs still operate within a single branch and do not explicitly compare different ways of reusing information across branches and graph-level summaries within one framework.

This paper studies Enhanced Cross Residual Graph Neural Networks (ECR-GNN), a modular framework for comparing several residual and cross-residual mechanisms within graph classification. Rather than treating the method as a single monolithic model, we formulate it as a family of six related graph classifiers built from a shared message-passing backbone. The family spans plain stacking, intra-branch residual reuse, node-level cross residual, sequential graph-level conditioning, repeated graph-level residual refinement, and graph-level cross residual between block pairs. This perspective makes it possible to ask a more practical question: under what graph conditions does extra cross-branch interaction help beyond a simpler residual baseline? In particular, we are interested in whether these mechanisms are most useful when graph labels depend on coordinated local structures and whole-graph arrangement rather than on a single dominant neighborhood pattern.

The contributions of this work are summarized as follows:

\begin{itemize}
    \item We formalize ECR-GNN as a unified family of graph classification architectures and distinguish the roles of node-level residual reuse and graph-level residual propagation within that family.
    \item We distinguish two reusable information-flow mechanisms: node-level cross residual between parallel branches and graph-level residual propagation between sequential blocks.
    \item We perform a controlled empirical study on four TUDataset benchmarks, covering six architectures, three operators, five depths, and two widths. The results show that residual reuse is the most reliable improvement, while cross-branch interaction is most effective on graph classification tasks that require stronger coupling between local substructures and global context.
    \item We provide a comparative analysis of accuracy, stability, and efficiency, together with a pseudocode-level description that highlights the common structure of the proposed architecture family.
\end{itemize}
